\clearpage
\thispagestyle{plain}
\markboth{CHAPITRE II}{CHAPITRE II}
\oldpage[II]{5}
\phantomsection
\label{chapter:II-fr}

\begin{center}
{\large CHAPITRE II\par}
\vspace{1.5em}
{\Large\bfseries ÉTUDE GLOBALE ÉLÉMENTAIRE\par}
{\Large\bfseries DE QUELQUES CLASSES DE MORPHISMES\par}
\vspace{1.8em}
{\bfseries Sommaire\par}
\end{center}

\medskip
\begin{tabular}{@{}r@{\ }p{0.82\textwidth}@{}}
\S~1. & Morphismes affines.\\
\S~2. & Spectres premiers homogènes.\\
\S~3. & Spectre homogène d'un faisceau d'algèbres graduées.\\
\S~4. & Fibrés projectifs. Faisceaux amples.\\
\S~5. & Morphismes quasi-affines~; morphismes quasi-projectifs~; morphismes
         propres~; morphismes projectifs.\\
\S~6. & Morphismes entiers et morphismes finis.\\
\S~7. & Critères valuatifs.\\
\S~8. & Schémas éclatés~; cônes projetants~; fermeture projective.
\end{tabular}

\medskip
Les diverses classes de morphismes étudiées dans ce chapitre le sont sans
faire grand usage des méthodes cohomologiques~; une étude plus poussée,
utilisant ces dernières méthodes, en sera faite au chapitre~III, où on
utilisera surtout les \S\S~2, 4 et~5 du chapitre~II. Le \S~8 peut être omis en
première lecture~: il donne quelques compléments au formalisme développé dans
les \S\S~1 à~3, se réduisant à des applications faciles de ce formalisme, et
nous en ferons un usage moins constant que des autres résultats de ce
chapitre.
